\documentclass[11pt]{article}
\usepackage[T1]{fontenc}
\usepackage[margin=1in]{geometry}
\usepackage{amsmath,amssymb,amsthm,mathtools}
\usepackage{graphicx}
\usepackage{microtype}
\usepackage[hidelinks]{hyperref}

\newtheorem{theorem}{Theorem}
\newtheorem{lemma}[theorem]{Lemma}
\newtheorem{remark}[theorem]{Remark}
\newcommand{\R}{\mathbb R}
\newcommand{\supp}{\operatorname{supp}}
\newcommand{\Mc}{\mathcal M_c}

\title{The General Eberlein--Grothendieck/\\Hemicompactness Equivalence}
\author{Candidate full solution for arXiv:2206.10684}
\date{13 August 2026}

\begin{document}
\maketitle

\begin{center}
\fbox{\parbox{0.91\linewidth}{\textbf{Status.}
Candidate full solution, likely valid; human review required.  The theorem
below removes the first-countability hypothesis from the result of
K\k{a}kol--Leiderman.  Novelty is provisional pending specialist review.}}
\end{center}

\section{Source problem}

K\k{a}kol and Leiderman study, for a Tychonoff space $X$, the locally convex
space $C_k(X)$ of real-valued continuous functions with the compact-open
topology.  They write $w$ for its weak topology and ask the following in
Problem~2.6 of \cite{source} (PDF p.~5):

\begin{quote}
Is $(C_k(X),w)$ Eberlein--Grothendieck if and only if $X$ is hemicompact?
\end{quote}

Here a space is Eberlein--Grothendieck if it embeds into $C_p(L)$ for some
compact $L$, and $X$ is \emph{hemicompact} if it has a sequence of compact
subsets cofinal among all compact subsets.  The source proves the equivalence
when $X$ is first countable.  Its exact statement and scope are shown below.

\begin{figure}[ht]
\centering
\includegraphics[width=\linewidth]{figures/open_problem_crop.png}
\caption{Problem~2.6 in \cite{source}, together with the source's
first-countable scope.}
\end{figure}

\section{Proof intuition}

The source reduces the Eberlein--Grothendieck condition to weak-star
$\sigma$-compactness of the dual.  Its dual elements are compactly supported
Radon measures.  The key new observation is that the supports of a weak-star
compact family cannot collectively escape along the level sets of a continuous
function: if they did, support-localized test functions in disjoint high
bands could be weighted and summed to form one continuous function on which
the compact family has unbounded evaluations.

Thus every compact piece of the dual has functionally bounded total support.
The Eberlein--Grothendieck hypothesis already forces $X$ itself to be
$\sigma$-compact.  Such a Tychonoff space is a paracompact $\mu$-space, so a
functionally bounded subset has compact closure.  An increasing compact cover
of the dual therefore gives increasing compact support envelopes $H_n$ in
$X$.

The remaining issue is cofinality, not coverage.  If a compact $K$ escaped
every $H_n$, choose $x_n\in K\setminus H_n$ and form the positive measure
$\sum 2^{-n}\delta_{x_n}$.  It belongs to some fixed dual-cover piece, forcing
all its atoms into that piece's support envelope.  Monotonicity says that all
sufficiently late atoms lie outside the same envelope, a contradiction.

\section{Preliminaries}

We use three results recorded in the source \cite{source}.

\begin{enumerate}
\item For a locally convex space $E$, $(E,w)$ is
Eberlein--Grothendieck exactly when $(E^*,w^*)$ is $\sigma$-compact
(Proposition~2.1).
\item Every $\phi\in C_k(X)^*$ is represented by a unique Radon measure
$\mu$ on the Hewitt realcompactification $\nu X$ whose support is a compact
subset of $X$ (Theorem~2.7, quoted from \cite{schmets}).  We identify
$C_k(X)^*$ with this measure space $\Mc(X)$.
\item The Dirac map $x\mapsto\delta_x$ is a homeomorphic closed embedding of
$X$ into the weak-star dual
\[
 (\Mc(X),\sigma(\Mc(X),C(X)))
\]
(Proposition~2.8).
\end{enumerate}

We also use the support-localization clause of the measure representation:
if $U\subset\nu X$ is open and meets $\supp\mu$, there is $g\in C(X)$ that
vanishes on $\nu X\setminus U$ and satisfies $\mu(g)=1$.

Recall that $B\subset X$ is \emph{functionally bounded} if every member of
$C(X)$ is bounded on $B$.

\begin{lemma}[Paracompact $\mu$-space fact]\label{lem:mu-space}
If $X$ is paracompact Hausdorff, then the closure of every functionally
bounded subset of $X$ is compact.
\end{lemma}

\begin{proof}
Let $B$ be functionally bounded and put $F=\overline B^X$.  If $F$ were
noncompact, then the paracompact Hausdorff space $F$ would not be
pseudocompact: a paracompact pseudocompact Hausdorff space is compact.
Consequently some $h\in C(F)$ would be unbounded.  The ambient space $X$ is
normal, and the real-valued Tietze extension theorem extends $h$ to a member
of $C(X)$.  Since $B$ is dense in $F$, continuity makes $h$ unbounded already
on $B$, contrary to functional boundedness.  Hence $F$ is compact.
\end{proof}

\clearpage
\section{The support-envelope lemma}

\begin{lemma}\label{lem:support}
Let $X$ be Tychonoff and let $A$ be a weak-star compact subset of $C_k(X)^*$.
Then
\[
 S_A:=\bigcup_{\mu\in A}\supp\mu
\]
is functionally bounded in $X$.
\end{lemma}

\begin{proof}
Suppose not.  Choose $u\in C(X)$ such that $|u|$ is unbounded on $S_A$, and
let $q=|\widetilde u|\in C(\nu X)$, where $\widetilde u$ is the extension of
$u$ to the Hewitt realcompactification.

We inductively choose measures $\mu_n\in A$, points
$x_n\in\supp\mu_n$, and pairwise disjoint open bands
\[
 U_n=\{y\in\nu X:a_n<q(y)<b_n\}.
\]
At stage $n$, all earlier supports are compact, so $q$ is bounded on their
finite union.  Since $q$ is unbounded on $S_A$, choose $x_n$ so high that the
interval $(a_n,b_n)$ containing $q(x_n)$ has its lower endpoint above the
previous bands, above the maximum of $q$ on every earlier support, and above
$n$.  Choose $\mu_n\in A$ with $x_n\in\supp\mu_n$.  At the next stage also
place the new lower endpoint above $\max_{\supp\mu_n}q$.  It follows that
\begin{equation}\label{eq:triangular}
 U_j\cap\supp\mu_i=\varnothing \qquad(j>i).
\end{equation}
The lower endpoints tend to infinity, so $(U_n)$ is locally finite.

By support localization, for every $n$ there is $g_n\in C(X)$ vanishing off
$U_n$ and satisfying $\mu_n(g_n)=1$.  Choose real numbers $c_n$ recursively.
After $c_1,\ldots,c_{n-1}$ are fixed, take
\[
 c_n=n-\sum_{j<n}c_j\mu_n(g_j).
\]
Because the bands are locally finite,
\[
 g:=\sum_{n=1}^{\infty}c_ng_n
\]
is a well-defined continuous function on $X$.  By
\eqref{eq:triangular}, every term $g_j$ with $j>n$ vanishes on
$\supp\mu_n$.  Hence
\[
 \mu_n(g)=\sum_{j\le n}c_j\mu_n(g_j)=n.
\]
On the other hand, evaluation at $g$ is a continuous real-valued function on
the weak-star compact space $A$, so it must be bounded.  This contradiction
proves the lemma.
\end{proof}

\begin{remark}
Lemma~\ref{lem:support} deliberately concludes only functional boundedness.
Without additional hypotheses on $X$, compactness of
$\overline{S_A}^{,X}$ can fail; this is why the $\sigma$-compactness forced
by the Eberlein--Grothendieck assumption is used separately below.
\end{remark}

\section{Full solution}

\begin{theorem}\label{thm:main}
For every Tychonoff space $X$, the following are equivalent:
\begin{enumerate}
\item $(C_k(X),w)$ is Eberlein--Grothendieck;
\item $X$ is hemicompact.
\end{enumerate}
\end{theorem}

\begin{proof}
Assume first that $(C_k(X),w)$ is Eberlein--Grothendieck.  By the source's
Proposition~2.1, the weak-star dual $\Mc(X)$ is $\sigma$-compact.  Choose
compact sets whose union is $\Mc(X)$ and replace them by their finite unions;
we obtain an increasing compact cover
\[
 A_1\subset A_2\subset\cdots,\qquad \Mc(X)=\bigcup_{n=1}^{\infty}A_n.
\]

The closed Dirac embedding of $X$ into $\Mc(X)$ shows that $X$ is
$\sigma$-compact.  A Tychonoff $\sigma$-compact space is Lindel\"of and hence
paracompact.  For every $n$, define
\[
 H_n=\overline{\bigcup_{\mu\in A_n}\supp\mu}^{,X}.
\]
Lemma~\ref{lem:support} and Lemma~\ref{lem:mu-space} show that each $H_n$ is
compact.  Moreover $(H_n)$ is increasing.

We claim that $(H_n)$ is cofinal among the compact subsets of $X$.  If not,
there is a compact $K\subset X$ with $K\nsubseteq H_n$ for every $n$.  Choose
\[
 x_n\in K\setminus H_n
\]
and consider the positive atomic measure
\[
 \mu=\sum_{n=1}^{\infty}2^{-n}\delta_{x_n}.
\]
This is a Radon probability supported in $K$.  For every $f\in C(X)$,
\[
 |\mu(f)|\leq \sup_{x\in K}|f(x)|,
\]
so $\mu$ is a continuous linear functional on $C_k(X)$, that is,
$\mu\in\Mc(X)$.  Choose $N$ with $\mu\in A_N$.  By definition of $H_N$,
\[
 \supp\mu\subset H_N.
\]
Every $x_n$ is a positive atom of $\mu$, so $x_n\in\supp\mu$.  But if
$n\geq N$, then $H_N\subset H_n$ and $x_n\notin H_n$, hence
$x_n\notin H_N$.  This is a contradiction.  Therefore $(H_n)$ is cofinal and
$X$ is hemicompact.

Conversely, if $X$ is hemicompact, then $C_k(X)$ is metrizable
\cite[Theorem~5.8.5]{narici}.  Every metrizable locally convex space endowed
with its weak topology is Eberlein--Grothendieck
\cite[Corollary~2.3]{source}.  Hence $(C_k(X),w)$ is
Eberlein--Grothendieck.
\end{proof}

\section{Verification, scope, and novelty}

The proof uses no first-countability, sequentiality, separability, or local
compactness.  The potentially delicate points are as follows.

\begin{itemize}
\item The support bands in Lemma~\ref{lem:support} are locally finite because
their lower levels tend to infinity; arbitrary coefficients therefore do not
endanger continuity of their sum.
\item The induction gives a triangular vanishing relation, so future
summands vanish on each earlier compact support.  There is no interchange of
an uncontrolled infinite sum and a signed measure.
\item The diagonal measure is positive.  Thus every chosen point is an atom
in its support; no compact set is assumed to carry a measure of full support.
\item Functional boundedness is upgraded to compact closure only after the EG
hypothesis forces $X$ to be $\sigma$-compact and paracompact.
\end{itemize}

The run's lightweight indexes contained no earlier result for arXiv:2206.10684.
Bounded exact-title, exact-problem, notation, author, hemicompactness, and
weak-star-support searches through 13 August 2026 found the source, a
conference presentation repeating the first-countable result, and later
citations, but no explicit resolution of Problem~2.6.  The argument should
therefore be treated as plausibly new, not as definitively novel, until an
expert literature review is completed.

\section{Human-review recommendation}

The candidate should be reviewed as a full solution.  The highest-value checks
are: (i) the support-localization theorem in the precise form used in
Lemma~\ref{lem:support}; (ii) the paracompact $\mu$-space fact; and (iii) the
support inclusion and monotonicity in the final atomic diagonal.  No
computational verification is relevant.

\begin{thebibliography}{9}
\raggedright

\bibitem{source}
J. K\k{a}kol and A. Leiderman,
\emph{When is a locally convex space Eberlein--Grothendieck?},
arXiv:2206.10684 (2022).

\bibitem{narici}
L. Narici and E. Beckenstein,
\emph{Topological Vector Spaces}, 2nd ed., Chapman and Hall/CRC, 2010.

\bibitem{schmets}
J. Schmets,
\emph{Espaces de fonctions continues},
Lecture Notes in Mathematics 519, Springer, 1976.

\end{thebibliography}

\end{document}
